\documentclass[12pt,a4paper]{article}
\usepackage[utf8]{inputenc}
\usepackage[T1]{fontenc}
\usepackage{lmodern}
\usepackage[brazil]{babel}
\usepackage[margin=2.5cm]{geometry}
\usepackage{amsmath,amssymb}
\usepackage{booktabs}
\usepackage{listings}
\usepackage{xcolor}
\usepackage{graphicx}
\usepackage{pgfplots}
\usepackage{hyperref}
\pgfplotsset{compat=1.16}
\hypersetup{colorlinks=true,linkcolor=black,citecolor=blue!50!black,urlcolor=blue!50!black}
\setlength{\parindent}{1.2em}
\setlength{\parskip}{0.3em}

\lstdefinestyle{cpp}{
  language=C++,
  basicstyle=\ttfamily\footnotesize,
  keywordstyle=\color{blue!60!black},
  commentstyle=\color{green!40!black}\itshape,
  numbers=left,
  numberstyle=\tiny\color{gray},
  breaklines=true,
  frame=single,
  framerule=0.3pt,
  columns=fullflexible,
  keepspaces=true,
  tabsize=2,
  extendedchars=true,
  literate=
    {á}{{\'a}}1 {à}{{\`a}}1 {â}{{\^a}}1 {ã}{{\~a}}1
    {é}{{\'e}}1 {ê}{{\^e}}1
    {í}{{\'i}}1
    {ó}{{\'o}}1 {ô}{{\^o}}1 {õ}{{\~o}}1
    {ú}{{\'u}}1 {ü}{{\"u}}1
    {ç}{{\c c}}1
    {Á}{{\'A}}1 {É}{{\'E}}1 {Í}{{\'I}}1 {Ó}{{\'O}}1 {Ú}{{\'U}}1 {Ç}{{\c C}}1
    {ε}{{$\epsilon$}}1 {α}{{$\alpha$}}1
}
\lstset{style=cpp}

\newcommand{\codigo}[3]{\lstinputlisting[caption={\texttt{#1}},label={lst:#2}]{../../trabalho/#3}}
\newcommand{\figura}[3]{\begin{center}\includegraphics[width=#2]{figuras/#1.pdf}\\[0.2em]\small #3\end{center}}
\newtheorem{teorema}{Teorema}
\newtheorem{lema}{Lema}
\newenvironment{demonstracao}{\par\noindent\textit{Demonstração.}\ }{\hfill$\square$\par}

\newcommand{\graficotempo}{
\begin{tikzpicture}\begin{semilogyaxis}[width=0.9\linewidth,height=6.5cm,xlabel={$n$},ylabel={Tempo mediano (ms)},legend pos=north west,grid=both,xtick={6,7,8,9,10,11,12,13,14},ymin=1e-3]
\addplot[color=red,mark=*] coordinates {(6,0.002) (7,0.009) (8,0.075) (9,0.751) (10,6.029) (11,57.604)};\addlegendentry{Força bruta}
\addplot[color=blue,mark=*] coordinates {(6,0.002) (7,0.010) (8,0.053) (9,0.380) (10,3.049) (11,38.012)};\addlegendentry{Backtracking}
\addplot[color=black,mark=*] coordinates {(6,0.001) (7,0.001) (8,0.003) (9,0.005) (10,0.026) (11,0.070) (12,0.150) (13,0.385) (14,0.955)};\addlegendentry{Held-Karp}
\end{semilogyaxis}\end{tikzpicture}}
\newcommand{\graficonos}{
\begin{tikzpicture}\begin{semilogyaxis}[width=0.9\linewidth,height=6.5cm,xlabel={$n$},ylabel={Mediana de nós},legend pos=north west,grid=both,xtick={6,7,8,9,10,11,12,13,14},ymin=1e-3]
\addplot[color=red,mark=*] coordinates {(6,120.000) (7,720.000) (8,5040.000) (9,40320.000) (10,362880.000) (11,3628800.000)};\addlegendentry{Força bruta}
\addplot[color=blue,mark=*] coordinates {(6,81.000) (7,391.000) (8,2283.000) (9,15657.000) (10,123301.000) (11,1096011.000)};\addlegendentry{Backtracking}
\addplot[color=black,mark=*] coordinates {(6,41.000) (7,97.000) (8,225.000) (9,513.000) (10,1153.000) (11,2561.000) (12,5633.000) (13,12289.000) (14,26625.000)};\addlegendentry{Held-Karp}
\end{semilogyaxis}\end{tikzpicture}}
\newcommand{\tabelafiltro}{
\begin{tabular}{r rr}\toprule $n$ & nós com filtro & nós sem filtro\\ \midrule
6 & 0 & 6\\
7 & 0 & 16\\
8 & 0 & 14\\
9 & 0 & 15\\
10 & 0 & 44\\
11 & 0 & 31\\
\bottomrule\end{tabular}}
\newcommand{\tabelapiorcaso}{
\begin{tabular}{r rr rr rr}\toprule & \multicolumn{2}{c}{Força bruta} & \multicolumn{2}{c}{Backtracking} & \multicolumn{2}{c}{Held-Karp}\\ $n$ & nós & ms & nós & ms & nós & ms\\ \midrule
6 & 120 & 0.002 & 81 & 0.002 & 41 & 0.001\\
7 & 720 & 0.009 & 391 & 0.010 & 97 & 0.001\\
8 & 5040 & 0.075 & 2283 & 0.053 & 225 & 0.003\\
9 & 40320 & 0.751 & 15657 & 0.380 & 513 & 0.005\\
10 & 362880 & 6.029 & 123301 & 3.049 & 1153 & 0.026\\
11 & 3628800 & 57.604 & 1096011 & 38.012 & 2561 & 0.070\\
12 & -- & -- & -- & -- & 5633 & 0.150\\
13 & -- & -- & -- & -- & 12289 & 0.385\\
14 & -- & -- & -- & -- & 26625 & 0.955\\
\bottomrule\end{tabular}}
\newcommand{\tabelaheuristicas}{
\begin{tabular}{l r rrr}\toprule Tipo & Digrafos com ciclo & Guloso índice & Guloso Warnsdorff & Recozimento\\ \midrule
Ciclo + arcos extras & 72 & 48/48 & 10/48 & 47/48\\
DAG aleatório & 0 & -- & -- & --\\
Erdős–Rényi, $p=0{,}35$ & 41 & 1/27 & 0/27 & 27/27\\
Erdős–Rényi, $p=\ln n/n$ & 3 & 0/3 & 0/3 & 3/3\\
Torneio aleatório & 62 & 6/42 & 3/42 & 42/42\\
Pior caso (sem ciclo) & 0 & -- & -- & --\\
\bottomrule\end{tabular}}
\newcommand{\tabelatipos}{
\begin{tabular}{l rr rr rr}\toprule & \multicolumn{2}{c}{Força bruta} & \multicolumn{2}{c}{Backtracking (Warnsdorff)} & \multicolumn{2}{c}{Held-Karp}\\ Tipo ($n=11$) & mediana & máximo & mediana & máximo & mediana & máximo\\ \midrule
Ciclo + arcos extras & 0.000 & 0.000 & 0.001 & 0.002 & 0.004 & 0.049\\
DAG aleatório & 0.000 & 0.000 & 0.000 & 0.000 & 0.000 & 0.001\\
Erdős–Rényi, $p=0{,}35$ & 3.131 & 36.111 & 0.005 & 0.030 & 0.016 & 0.044\\
Erdős–Rényi, $p=\ln n/n$ & 0.000 & 0.000 & 0.000 & 0.000 & 0.001 & 0.001\\
Torneio aleatório & 0.822 & 14.243 & 0.002 & 0.039 & 0.058 & 0.076\\
\bottomrule\end{tabular}}
\newcommand{\totalmedicoes}{2528}\newcommand{\totalcensurados}{0}


\title{\textbf{Implementação e Análise de Estratégias de Algoritmos\\para Solução de Problema NP}\\[0.6em]
\large Ciclo Hamiltoniano Direcionado}
\author{}
\date{}

\begin{document}
\maketitle
\thispagestyle{empty}

\section*{Autores}

\begin{center}
\begin{tabular}{ll}
\toprule
Nome completo & Matrícula\\
\midrule
Fabio Firanzi & \textbf{[matrícula]}\\
Finéias Antonio & \textbf{[matrícula]}\\
Heitor Vaz & \textbf{[matrícula]}\\
Tulio Vitorette & \textbf{[matrícula]}\\
\bottomrule
\end{tabular}
\end{center}

\tableofcontents
\newpage

\section{Introdução}

\subsection{O problema}

Seja $G=(V,A)$ um grafo direcionado simples, sem laços, com $n=|V|\ge 2$ vértices. Um \emph{ciclo hamiltoniano direcionado} é uma sequência $(v_1,v_2,\dots,v_n)$ tal que

\begin{itemize}
  \item os $v_i$ são todos distintos e cobrem $V$;
  \item $(v_i,v_{i+1})\in A$ para $i=1,\dots,n-1$;
  \item $(v_n,v_1)\in A$.
\end{itemize}

O \textbf{problema de decisão (CHD)} pergunta: dado $G$, existe um ciclo hamiltoniano direcionado? O código deste trabalho resolve a \emph{versão de busca}: devolve a sequência quando ela existe. A resposta de decisão é o sim ou o não dessa busca, e a própria sequência é a prova de que a resposta é sim, usada na Seção~\ref{sec:prova}.

\paragraph{Convenção de simetria.} Rotacionar um ciclo não cria outro ciclo. Por isso todos os métodos exatos fixam $v_1=0$. Inverter o sentido, ao contrário, produz outro ciclo direcionado e não é equivalente: o sentido oposto continua valendo e entra na busca.

\subsection{Dois exemplos mínimos}

\paragraph{Exemplo com ciclo.} Arcos $(0,1),(1,2),(2,3),(3,0)$. O ciclo é $0\to1\to2\to3\to0$, em verde na figura.

\figura{ex_sim}{0.42\linewidth}{Digrafo com ciclo hamiltoniano direcionado.}

\paragraph{Exemplo sem ciclo, com o digrafo forte.} Arcos $(0,1),(1,2),(2,0),(1,3),(3,0)$. De $3$ só se sai para $0$, e em $3$ só se entra por $1$, então $3$ teria de ocupar o trecho $1\to3\to0$. O vértice $2$ não cabe nesse trecho com os arcos existentes. O digrafo é fortemente conexo: conectividade forte é condição \emph{necessária} e não suficiente.

\figura{ex_nao}{0.5\linewidth}{Digrafo fortemente conexo sem ciclo hamiltoniano.}

Os dois exemplos podem ser executados por qualquer programa do trabalho, com os tipos \texttt{ciclo\_simples} e \texttt{forte\_sem\_ciclo}.

\subsection{Aplicações práticas}
\label{sec:aplicacoes}

Em cada aplicação abaixo, o que é vértice, o que é arco e por que a volta ao início importa. A dificuldade é a mesma nos quatro casos: o ciclo hamiltoniano é NP-completo, e a solução exata só funciona para digrafos pequenos.

\paragraph{Entregas urbanas (Correios, iFood, Amazon).} Cada endereço de entrega é um vértice, e um arco é o trecho de rua que leva de um endereço a outro. Como há ruas de mão única, ir de $A$ a $B$ não equivale a ir de $B$ a $A$, e o digrafo é naturalmente \emph{direcionado}. O entregador parte do depósito, visita cada endereço uma vez e volta ao depósito: a volta ao início é o retorno da van ou da moto à base. As empresas resolvem a versão com custos (caixeiro viajante assimétrico e roteirização de veículos), e o CHD é a pergunta de viabilidade dentro dela: o caso em que todo arco existente custa $1$ e o limite é $B=n$. Com centenas de paradas por dia, a solução exata é inviável (Seção~\ref{sec:complexidade}) e as empresas usam heurísticas, como o recozimento deste trabalho. O texto não atribui a elas razão de aproximação provada.

\paragraph{Computação com DNA (Adleman, 1994).} A primeira demonstração de computação com moléculas de DNA resolveu um problema de caminho hamiltoniano num \emph{digrafo} de $7$ vértices \cite{adleman1994}. Cada cidade foi codificada por uma sequência de DNA e cada arco por uma sequência que liga o fim de uma cidade ao início da seguinte. Misturadas num tubo de ensaio, as moléculas formaram, em paralelo, todos os caminhos possíveis, e as etapas seguintes eliminaram os que não passavam por todas as cidades exatamente uma vez. O experimento é famoso por mostrar a ideia, e também por mostrar o limite: o número de moléculas necessárias cresce com o número de caminhos, isto é, de forma exponencial, e a abordagem não escala.

\paragraph{Montagem de genomas (Projeto Genoma Humano e bactérias).} Para sequenciar um genoma, a máquina lê milhões de fragmentos curtos. Na formulação clássica, cada fragmento é um vértice, e um arco liga dois fragmentos quando o final de um se sobrepõe ao começo do outro. Uma sequência que usa cada fragmento uma vez é um caminho hamiltoniano; no genoma circular de uma bactéria ou de um plasmídeo é um ciclo hamiltoniano, e a volta ao início é a própria circularidade do DNA. Como o problema é NP-completo, os montadores modernos mudaram de modelo: usam o grafo de de Bruijn e o caminho euleriano, que é polinomial \cite{compeau2011}. O problema deste trabalho explica uma mudança real de abordagem numa área muito conhecida.

\paragraph{Campeonato por pontos corridos (Brasileirão, xadrez).} Num campeonato em que todos jogam contra todos, os competidores são vértices e uma vitória é um arco do vencedor para o perdedor. Todo par tem exatamente um sentido (sem empate, ou desconsiderando-o), então o digrafo é um torneio. Se ele é fortemente conexo, o teorema de Camion \cite{camion1959} garante um ciclo hamiltoniano: uma ordenação circular em que cada competidor venceu o seguinte, e o último venceu o primeiro. O ciclo explica por que o ``ganhou de'' não é transitivo (A vence B, B vence C, C vence A). Este é o caso em que a pergunta deixa de ser NP-completa: a construção da Seção~\ref{sec:camion} acha o ciclo em tempo polinomial, e mostra que a dificuldade mora na estrutura do digrafo e não no nome do problema.

\section{Prova: CHD pertence à classe NP}
\label{sec:prova}

\subsection{CHD $\in$ NP}

Um problema está em NP quando, se alguém propõe uma resposta ``sim'', dá para conferi-la rapidamente, em tempo polinomial \cite{cormen2009,ziviani2006}. Não é preciso saber \emph{achar} a resposta, só \emph{conferir}.

No CHD, a resposta ``sim'' proposta é uma sequência de vértices, $c=(c_1,\dots,c_n)$. Para conferir se ela é um ciclo hamiltoniano direcionado (função \texttt{verifica\_ciclo}), basta checar três coisas:

\begin{enumerate}
  \item A sequência tem exatamente $n$ vértices.
  \item Nenhum vértice se repete. Um vetor de $n$ marcas resolve em $n$ passos. Como são $n$ vértices diferentes em $n$ posições, nenhum vértice fica de fora.
  \item Cada vértice tem arco para o próximo, e o último tem arco para o primeiro. São $n$ consultas à matriz, cada uma em tempo constante.
\end{enumerate}

Ao todo são cerca de $n$ passos, um tempo polinomial. A conferência aceita se, e somente se, $c$ é de fato um ciclo hamiltoniano direcionado de $G$. Logo, CHD está em NP.

\subsection{CHD é NP-completo (além do que o edital pede)}

Um problema é NP-completo quando está em NP e é pelo menos tão difícil quanto os problemas mais difíceis de NP. Para mostrar isso, usa-se um problema que já se sabe ser difícil: o ciclo hamiltoniano em grafos \emph{sem direção}, que é NP-completo \cite{karp1972,garey1979}.

A ideia é transformar um problema no outro:

\begin{enumerate}
  \item Pegue um grafo sem direção e troque cada aresta entre $u$ e $v$ por dois arcos, um de $u$ para $v$ e outro de $v$ para $u$. Isso gera um digrafo.
  \item Se o grafo tem um ciclo hamiltoniano, o digrafo tem o mesmo ciclo, pois todo caminho existe nos dois sentidos.
  \item Se o digrafo tem um ciclo hamiltoniano, o grafo também tem, pois cada arco usado veio de uma aresta.
\end{enumerate}

Logo, quem soubesse resolver o CHD saberia resolver o ciclo hamiltoniano sem direção, que é um problema difícil. Então o CHD também é difícil e, como já está em NP, é NP-completo (problema GT38 de \cite{garey1979}).

\section{Implementação}

\subsection{Estrutura comum das implementações}

Cada estratégia é um programa independente, com a sua própria função \texttt{main}. Todos usam as mesmas peças, descritas abaixo, para que as comparações sejam justas: o digrafo, o contador e o verificador são iguais para todos.

\paragraph{Organização dos arquivos.} Há uma pasta para a parte comum e uma para cada estratégia. Cada pasta de estratégia tem o arquivo com o algoritmo (por exemplo, \texttt{forca\_bruta.cpp}) e o seu \texttt{main.cpp}, que apenas prepara o digrafo, chama o algoritmo e mostra o resultado. O comando \texttt{make} compila todos os programas.

\paragraph{O digrafo.} É guardado como uma tabela $n\times n$, em que a posição $(u,v)$ indica se existe o arco de $u$ para $v$. Assim, perguntar se um arco existe custa um passo, o que importa porque todos os algoritmos fazem essa pergunta muitas vezes. A tabela ocupa $n^2$ posições. Não existem laços (arcos de um vértice para ele mesmo): se alguém tenta adicionar um, ele é ignorado. A tabela distingue os sentidos, então a existência de $u\to v$ não implica a de $v\to u$.

\paragraph{O contador de operações.} Para a análise por contagem de instruções, o digrafo conta toda consulta de arco que recebe, e cada algoritmo conta também a sua unidade de trabalho principal. São três contadores:
\begin{itemize}
  \item \texttt{consultas\_arco}: quantas vezes se perguntou se um arco existe. Todos os algoritmos usam.
  \item \texttt{permutacoes}: quantas ordens de vértices foram testadas. Só a força bruta usa.
  \item \texttt{nos}: a unidade de trabalho dos demais algoritmos, que muda de um para outro: chamadas da busca no backtracking, estados expandidos no Held-Karp, escolhas no guloso e movimentos no recozimento.
\end{itemize}
Os números de algoritmos diferentes não são somados em um total único, porque um estado do Held-Karp e uma permutação da força bruta não são a mesma operação.

\paragraph{O verificador.} A função \texttt{verifica\_ciclo} recebe um digrafo e uma sequência de vértices e responde se ela é um ciclo hamiltoniano direcionado: confere se tem $n$ vértices, se nenhum se repete e se cada vértice tem arco para o seguinte, inclusive do último para o primeiro. É o verificador da prova de NP da seção anterior. Todos os programas o usam para conferir o ciclo que o algoritmo devolveu, antes de mostrá-lo.

\paragraph{O gerador de digrafos.} Cria os digrafos de teste a partir de um tipo e do tamanho $n$:
\begin{center}
\small
\begin{tabular}{p{4.7cm} p{8.6cm}}
\toprule
Tipo & O que gera\\
\midrule
\texttt{ciclo\_simples} & Só o ciclo $0\to1\to\cdots\to(n-1)\to0$. Tem ciclo hamiltoniano, e é o único.\\
\texttt{forte\_sem\_ciclo} & O exemplo fixo de 4 vértices da introdução: fortemente conexo e sem ciclo hamiltoniano. Ignora o $n$.\\
\texttt{pior\_caso} & Um digrafo completo sobre $n-1$ vértices, mais um vértice ligado só ao vértice 1. Não tem ciclo, e obriga os métodos exatos a esgotar a busca.\\
\texttt{aleatorio} & Cada arco existe com probabilidade $0{,}35$. Pode ter ou não ter ciclo, e a chance de ter cresce com $n$.\\
\texttt{armadilha\_indice} & Digrafo de 4 vértices em que o guloso por índice falha, embora exista ciclo.\\
\texttt{armadilha\_warnsdorff} & Digrafo de 4 vértices em que o guloso de Warnsdorff falha, embora exista ciclo.\\
\bottomrule
\end{tabular}
\end{center}
O tipo \texttt{aleatorio} usa uma semente. Sem indicação, ela é sorteada a cada execução, e o programa imprime a semente usada. Com a opção \texttt{-{}-semente=N} o mesmo digrafo é gerado de novo, o que permite rodar dois algoritmos no mesmo digrafo e comparar as respostas.

\paragraph{O programa de cada estratégia.} Todos funcionam do mesmo jeito: recebem o tipo e o tamanho $n$ pela linha de comando, geram o digrafo, rodam o algoritmo e imprimem o ciclo (ou a informação de que não foi encontrado), os contadores, o tempo e a semente do digrafo. O tempo é medido com \texttt{std::chrono::steady\_clock} e cobre só a execução do algoritmo, sem a geração do digrafo.

\subsection{Força bruta}

Fixa o vértice $0$ na origem e gera as $(n-1)!$ permutações dos demais com \texttt{std::next\_permutation}. Para cada permutação consulta os $n$ arcos do ciclo e para na primeira que fecha. Não há poda: poda é o que distingue o backtracking na comparação.

\subsection{Backtracking}

Busca em profundidade a partir de $0$. O próximo vértice é um vizinho de saída ainda fora do caminho; com $n$ vértices no caminho, o único teste é o arco de volta a $0$. Duas ordens de vizinhos (mesma resposta, árvores de tamanhos diferentes): índice crescente e a regra de Warnsdorff adaptada (primeiro o vizinho com menos saídas livres, ideia herdada do passeio do cavalo, sem herdar uma prova). A ordenação é detalhe do backtracking, não um algoritmo guloso. O parâmetro \texttt{usar\_filtro} permite a ablação da poda por conectividade forte.

\subsection{Divisão e conquista: por que não há algoritmo}
\label{sec:dc}

Divisão e conquista, no sentido de \cite{cormen2009}, parte a entrada em subproblemas \emph{independentes}, resolve cada um e combina as respostas. No CHD a combinação não é independente: um ciclo hamiltoniano de $G$ não é a união de ciclos hamiltonianos dos subdigrafos induzidos por uma partição $V=L\cup R$, porque o ciclo pode alternar entre $L$ e $R$ várias vezes.

O atalho que \emph{parece} divisão (permutar um bloco $L$, permutar o outro bloco $R$ e concatenar) responde outra pergunta. Contraexemplo, com origem $0$, $L=\{1,2\}$ e $R=\{3\}$: as concatenações possíveis são só $(0,1,2,3)$, $(0,2,1,3)$, $(0,3,1,2)$ e $(0,3,2,1)$. O digrafo de arcos $(0,1),(1,3),(3,2),(2,0)$ tem o ciclo $0\to1\to3\to2\to0$ e nenhuma das quatro sequências, então o método de blocos responde \emph{não} num digrafo que \emph{tem} ciclo.

\figura{blocos}{0.42\linewidth}{O ciclo $0\to1\to3\to2\to0$ intercala os blocos $L=\{1,2\}$ (azul) e $R=\{3\}$ (laranja), e nenhuma concatenação de blocos o produz.}

Por isso o grupo \textbf{não apresenta} esse atalho como solução. Ele existiu apenas para demonstrar o falso negativo, que o programa \texttt{ciclo\_por\_blocos} exibe com o tipo \texttt{ciclo\_intercalado}.

A decomposição que funciona usa outro subproblema: ``caminho simples que começa em $0$, cobre o conjunto $S$ e termina em $v$''. Esses subproblemas se repetem ao longo da busca; resolver cada um uma vez e reutilizar é programação dinâmica, a próxima subseção. Algoritmos mais recentes (por exemplo, a inclusão-exclusão de Björklund para o caso não direcionado \cite{bjorklund2010}) ficam fora do código, e o algoritmo exato de referência do trabalho é Held-Karp.

\subsection{Programação dinâmica (Held-Karp)}

Estado: $\mathrm{dp}[S][v]$ é verdadeiro quando existe caminho simples que começa em $0$, visita exatamente o conjunto $S$ e termina em $v$, com $0\in S$ e $v\in S$. Base: $\mathrm{dp}[\{0\}][0]=$ verdadeiro. Transição: $\mathrm{dp}[S][v]$ é verdadeiro quando existe $u\in S\setminus\{v\}$ com $\mathrm{dp}[S\setminus\{v\}][u]$ verdadeiro e arco $(u,v)$. Há ciclo quando existe $v\ne0$ com $\mathrm{dp}[V][v]$ verdadeiro e arco $(v,0)$ \cite{held1962,bellman1962}.

Cada conjunto é uma máscara de bits, processada em ordem crescente (a máscara sem o bit de $v$ é menor e já foi calculada). Para cada máscara guarda-se um conjunto de vértices finais possíveis (um inteiro de 32 bits) e, para cada estado, o predecessor, de modo que o ciclo é reconstruído, e não só o booleano. A função aceita até $n=22$ e, acima disso, lança erro em vez de responder que o ciclo não existe.

\subsection{Guloso}
\label{sec:guloso}

Dois critérios, mesma estrutura: parte de $0$ e, enquanto houver vértice livre, escolhe um vizinho de saída livre; se empaca antes de cobrir $V$, ou se o último não aponta para $0$, devolve falha. O critério \texttt{indice} escolhe o menor índice; o critério \texttt{warnsdorff} escolhe o vizinho com menos saídas livres. \textbf{Nenhum dos dois tem garantia}, e cada um tem contraexemplo automatizado:

\paragraph{Critério \texttt{indice}.} Arcos $(0,1),(0,2),(2,3),(3,1),(1,0)$. O ciclo $0\to2\to3\to1\to0$ existe (verde). O guloso sai de $0$ para $1$ (vermelho), e de $1$ só há arco para $0$, que já está no caminho: falha num digrafo que tem ciclo.

\figura{guloso_indice}{0.42\linewidth}{Armadilha do critério por índice. Verde: ciclo existente. Vermelho: escolha do guloso.}

\paragraph{Critério \texttt{warnsdorff}.} Arcos $(0,1),(0,2),(1,2),(1,3),(2,3),(3,0)$. O ciclo $0\to1\to2\to3\to0$ existe (verde). Saindo de $0$, o vértice $1$ tem duas saídas livres e o $2$ tem uma; a regra escolhe $2$ (vermelho), segue a $3$ e empaca com $1$ de fora.

\figura{guloso_warns}{0.42\linewidth}{Armadilha do critério de Warnsdorff. Verde: ciclo existente. Vermelho: escolha do guloso. Tracejado: arco não usado.}

Os dois são heurísticas gulosas, sem garantia de otimalidade e sem razão de aproximação.

\subsection{Heurística: recozimento simulado}

Reformulação de otimização usada pela heurística: entre as permutações com origem $0$, maximizar o número de arcos presentes nas $n$ posições do ciclo. O \emph{fitness} vai de $0$ a $n$, e o digrafo tem ciclo exatamente quando o máximo vale $n$. Essa reformulação permite busca local; \textbf{ela não cria razão de aproximação}.

Movimento: retirar um vértice e inseri-lo em outra posição (relocação). A \emph{variação} do fitness depende só dos arcos da fronteira da posição antiga e da nova, e é calculada em $O(1)$. Aceita melhora sempre e aceita piora com probabilidade $e^{\Delta/T}$; esfriamento $T\leftarrow\alpha T$; vários recomeços a partir de permutações aleatórias; para ao achar fitness $n$ ou ao gastar o orçamento \cite{kirkpatrick1983}.

Parâmetros declarados antes da bateria e congelados: $T_0=1$, $\alpha=0{,}995$, $200\cdot n$ movimentos por recomeço e $10$ recomeços. O resultado separa \texttt{Ciclo} (ciclo válido encontrado) de \texttt{Orcamento} (orçamento esgotado). Esgotar o orçamento \textbf{não} prova que o ciclo não existe.

\subsection{Caso especial: torneios (Camion)}
\label{sec:camion}

\begin{teorema}[Camion, 1959 \cite{camion1959}]
Um torneio com $n\ge3$ vértices tem ciclo hamiltoniano direcionado se e somente se é fortemente conexo.
\end{teorema}

Este módulo não é uma das seis estratégias do caso geral: mostra a fronteira da dificuldade, pois a mesma pergunta, restrita a torneios fortes, cai em P. Construção: (i) caminho hamiltoniano por inserção; (ii) um arco de volta ao início do caminho fecha um ciclo que cobre um prefixo; (iii) cada vértice de fora que encaixa entre dois consecutivos do ciclo entra ali; (iv) quando nenhum encaixa sozinho, os vértices de fora se dividem entre os que apontam para o ciclo inteiro e os que recebem arco de todo o ciclo, e um arco do segundo grupo para o primeiro permite inserir um par. O módulo recusa o digrafo que não é torneio.

\section{Análise de complexidade}
\label{sec:complexidade}

\subsection{Tempo e espaço teóricos}

Usa-se $n=|V|$ e matriz de adjacência, com consulta de arco em $O(1)$.

\begin{center}
\small
\begin{tabular}{p{3.2cm} p{4.4cm} p{3.3cm} p{2.8cm}}
\toprule
Função & Tempo & Espaço extra & Garantia\\
\midrule
\texttt{verifica\_ciclo} & $O(n)$ & $O(n)$ & exata\\
Força bruta & $\Theta(n\cdot(n-1)!)$ no pior caso & $O(n)$ & exata\\
Backtracking & $O(n!)$ no pior caso; depende da densidade & $O(n)$ (pilha) & exata\\
Held-Karp & $\Theta(n^2\,2^n)$ & $\Theta(n\,2^n)$ & exata\\
Guloso, índice & $O(n^2\log n)$ & $O(n)$ & nenhuma\\
Guloso, Warnsdorff & $O(n^3)$ & $O(n)$ & nenhuma\\
Recozimento & $O(r\cdot c\,n^2)$, $r$ recomeços, $c=200$ & $O(n)$ & nenhuma\\
Camion (só torneio) & $O(n^3)$ & $O(n)$ & exata\\
Blocos contíguos & $O(n!\cdot n)$ ($n\le10$) & $O(n)$ & incorreto\\
\bottomrule
\end{tabular}
\end{center}

\paragraph{Força bruta.} Gera $(n-1)!$ permutações; cada uma consulta até $n$ arcos e copia $n$ posições. O melhor caso termina na primeira permutação; o pior caso é $\Theta(n\cdot(n-1)!)$, e é o caso dos digrafos sem ciclo. Espaço $O(n)$, pois \texttt{next\_permutation} gera no lugar.

\paragraph{Backtracking.} A árvore tem profundidade $n$ e grau no máximo $n-1$, portanto $O(n!)$ nós no pior caso. Com a poda de vértice já usado, o número de nós em um digrafo completo sobre $m$ vértices é $\sum_{k}m!/(m-k)!\approx e\cdot m!$. Para o espaço, a pilha de recursão e o vetor de usados são $O(n)$.

\paragraph{Held-Karp.} São $2^n$ máscaras, para cada uma até $n$ vértices finais e para cada vértice até $n$ vizinhos de saída: $\Theta(n^2\,2^n)$. O tempo depende de $n$ e não do número de arcos, exceto pelas podas por máscara inalcançável. O espaço é dominado pela tabela de predecessores, com $n\,2^n$ inteiros de 4 bytes, mais $2^n$ palavras de 32 bits do conjunto de finais. Para $n=22$ são $22\cdot2^{22}\cdot4\approx 369$~MB mais $16$~MB, aproximadamente $385$~MB. Por isso a memória limita o $n$ \emph{antes} do tempo, e o código recusa $n>22$.

\paragraph{Guloso.} Cada passo percorre os vizinhos de saída ($O(n)$) e ordena os candidatos ($O(n\log n)$); são $n-1$ passos, daí $O(n^2\log n)$. No critério de Warnsdorff, cada candidato conta suas próprias saídas livres ($O(n)$), o que leva a $O(n^3)$. (Trocar a ordenação por uma busca do mínimo reduziria para $O(n^2)$ e $O(n^3)$ sem mudar a resposta; não foi feito porque o objetivo do trabalho é a comparação.)

\paragraph{Recozimento.} O cálculo da variação é $O(1)$ (seis consultas de arco), mas aplicar a relocação aceita desloca até $n$ posições, então um movimento custa $O(n)$ quando aceito. Com $r$ recomeços de $c\,n$ movimentos, o tempo é $O(r\,c\,n^2)$, proporcional ao orçamento declarado. Espaço $O(n)$.

\paragraph{Camion.} O caminho por inserção custa $O(n^2)$ (cada vértice varre o caminho). Cada rodada de inserção percorre os vértices de fora e o ciclo, $O(n^2)$, e pode haver até $n$ rodadas, portanto $O(n^3)$ no pior caso. Para $n\le 30$ no experimento o tempo é irrelevante frente aos métodos exponenciais.

\subsection{O que o contador mede}

A contagem de instruções usa as três grandezas do \texttt{Contador}. A \emph{unidade comum} de comparação entre estratégias é o tempo de parede. Cada algoritmo conta ``nós'' de uma natureza distinta, e por isso o texto não soma contagens de algoritmos diferentes num ``total de instruções'' único.

\begin{center}
\small
\begin{tabular}{lp{9cm}}
\toprule
Algoritmo & O que é um ``nó''\\
\midrule
Força bruta & uma permutação testada (valor medido: exatamente $(n-1)!$ no pior caso)\\
Backtracking & uma chamada recursiva da busca\\
Held-Karp & um estado (máscara, vértice final) expandido\\
Guloso & uma escolha de vizinho\\
Recozimento & um movimento de relocação proposto\\
Camion & uma inserção\\
\bottomrule
\end{tabular}
\end{center}

\section{Análise de resultados}

\subsection{Método}

\paragraph{Digrafos de teste.} Semente base $20261107$. Seis tipos de digrafos, cada um com gerador próprio:
\begin{itemize}
  \item \textbf{Ciclo + arcos extras}: um ciclo plantado mais $n$ arcos aleatórios (resposta sim).
  \item \textbf{DAG aleatório}: arcos só para a frente, $p=0{,}5$ (resposta não; não é fortemente conexo).
  \item \textbf{Erdős--Rényi direcionado}, com $p=0{,}35$ fixo e com $p=(\ln n)/n$, a ordem do limiar de aparecimento de grau mínimo de entrada e saída $1$ \cite{frieze1988}. Uma amostra finita não é tratada como demonstração do teorema.
  \item \textbf{Torneio aleatório}: sim quando é forte.
  \item \textbf{Pior caso (sem ciclo)}: um digrafo completo sobre $n-1$ vértices mais um vértice $x$ ligado apenas ao vértice $1$ (arcos $1\to x$ e $x\to1$). O digrafo é fortemente conexo, não tem ciclo hamiltoniano (passar por $x$ exige usar $1$ duas vezes) e obriga os métodos exatos a esgotar a busca.
\end{itemize}

\paragraph{Juiz.} Até $n=14$ o Held-Karp julga os demais; até $n=11$ a força bruta concorda com ele (nos $288$ digrafos em que mais de um método exato rodou, as respostas coincidiram). Todo ciclo devolvido passa por \texttt{verifica\_ciclo}. Acima do alcance da PD, ``ciclo válido encontrado'' ou ``orçamento esgotado''; a palavra ``inexistente'' fica reservada aos métodos exatos.

\paragraph{Execução.} AMD Ryzen 7 5700G, 31~GiB de memória, g++ 16.2 com \texttt{-O2}, um núcleo. Foram \totalmedicoes{} medições, com $8$ digrafos por tamanho e tipo. Timeout de $30$~s por chamada exata: \totalcensurados{} medições foram censuradas. Tamanhos: força bruta e backtracking $n=6$ a $11$; Held-Karp $n=6$ a $14$; guloso, recozimento e Camion até $n=30$. Os números de nós são reproduzíveis com os programas do trabalho, pois a semente é fixa; os tempos variam um pouco de uma execução para outra.

\paragraph{Cuidado com a resolução do relógio.} Tempos abaixo de aproximadamente $0{,}005$~ms são ruído de medição (resolução de $0{,}001$~ms nos dados) e não são interpretados. A discussão se apoia em tempos acima de $0{,}01$~ms e nas contagens de nós, que não têm ruído.

\subsection{Os três exatos no pior caso}

\begin{center}
\tabelapiorcaso
\end{center}

A Figura~\ref{fig:tempo} e a Figura~\ref{fig:nos} mostram os mesmos dados em escala logarítmica.

\begin{figure}[h]
\centering
\graficotempo
\caption{Tempo mediano contra $n$ no tipo pior caso (escala logarítmica).}
\label{fig:tempo}
\end{figure}

\begin{figure}[h]
\centering
\graficonos
\caption{Mediana de nós contra $n$ no tipo pior caso (escala logarítmica).}
\label{fig:nos}
\end{figure}

\paragraph{Força bruta.} Os nós medidos são exatamente $(n-1)!$: $120$ em $n=6$ e $3\,628\,800$ em $n=11$, igual à teoria. O tempo cresce por um fator de cerca de $9{,}6$ entre $n=10$ e $n=11$ ($6{,}03\to57{,}6$~ms), compatível com o fator $n-1=10$.

\paragraph{Backtracking.} A poda de vértice usado reduz os nós a cerca de $30\%$ da força bruta em $n=11$ ($1\,096\,011$ contra $3\,628\,800$), valor próximo de $e\cdot9!\approx 986\,000$ da fórmula da seção anterior. O ganho é de uma constante, não da ordem: a razão entre nós consecutivos continua a crescer com $n$. Neste tipo as duas ordens de vizinhos têm o mesmo número de nós (esperado: o grafo é completo sobre $n-1$ vértices e o digrafo não tem ciclo, logo a árvore é esgotada de qualquer modo); a ordem de Warnsdorff custa mais tempo pelo cálculo da chave ($43{,}6$ contra $38{,}0$~ms).

\paragraph{Held-Karp.} Os nós crescem por um fator de aproximadamente $2{,}2$ por unidade de $n$ ($2561\to5633\to12289\to26625$), compatível com $\Theta(n^2\,2^n)$, e o tempo cresce por cerca de $2{,}5$ por unidade de $n$ ($0{,}150\to0{,}385\to0{,}955$~ms). Em $n=11$, Held-Karp leva $0{,}070$~ms contra $57{,}6$~ms da força bruta, uma razão de cerca de $800$ vezes, e a razão aumenta com $n$. O cruzamento entre o fatorial e $n^2\,2^n$ acontece já em $n\approx 7$, onde as duas curvas são indistinguíveis dentro do ruído; a partir daí a diferença é de ordem.

\subsection{Efeito do filtro de conectividade forte}

No tipo DAG a resposta é sempre não. Com o filtro, o backtracking responde sem expandir nenhum nó; sem o filtro, expande alguns:

\begin{center}
\tabelafiltro
\end{center}

O filtro corta a árvore inteira, mas é \emph{barato e correto}, e não resolve o caso forte: no tipo pior caso, que passa no filtro, a busca esgota $1{,}1$ milhão de nós em $n=11$.

\subsection{Variação entre tipos}

\begin{center}
\tabelatipos
\end{center}
\noindent Tempos em milissegundos, $n=11$, oito digrafos por tipo.

Nos tipos em que existe ciclo e há muitos arcos, o backtracking acha rapidamente (cauda curta). A força bruta tem cauda longa nos tipos com ciclo: em Erdős--Rényi com $p=0{,}35$ a mediana é de $3{,}1$~ms e o máximo, $36$~ms, pois a primeira permutação que fecha pode estar tarde na ordem lexicográfica. O Held-Karp é estável (o tempo muda pouco entre tipos do mesmo $n$ nos que existe ciclo), como esperado de uma complexidade que não depende da densidade. Mediana e máximo contam histórias diferentes, e os dois são publicados.

\subsection{Guloso e recozimento: quem acha o ciclo?}

\begin{center}
\tabelaheuristicas
\end{center}
\noindent Digrafos em que o Held-Karp (juiz) confirma ciclo, $n=6$ a $14$. A segunda coluna é o total de digrafos com ciclo; as heurísticas rodam só em parte dos tamanhos, então cada célula é encontrados / digrafos com ciclo \emph{em que a heurística rodou}.

\begin{itemize}
  \item \textbf{O guloso por índice acerta 48 de 48 no tipo ciclo + arcos extras por um efeito do gerador}: o ciclo plantado é $0\to1\to2\to\cdots$, que é exatamente a ordem de índice. Isso mede o viés do tipo, e não a qualidade do critério; nos demais tipos (Erdős--Rényi, torneio) o mesmo critério encontra $1/27$ e $6/42$.
  \item \textbf{A regra de Warnsdorff} é pior que o índice nesses digrafos ($10/48$, $0/27$, $3/42$). O contraexemplo da Seção~\ref{sec:guloso} mostra que é falha de projeto, e não azar.
  \item \textbf{O recozimento} encontra $47/48$, $27/27$, $3/3$ e $42/42$ dos ciclos, recuperando a maior parte das falhas dos gulosos, no mesmo orçamento fixo. Mas continua sem como provar que o ciclo não existe: nos digrafos sem ciclo (DAG, pior caso) ele sempre termina por \texttt{Orcamento}, e isso não prova nada.
  \item O tipo de limiar $p=(\ln n)/n$ tem só $3$ digrafos com ciclo no total. A amostra é pequena e não sustenta conclusão estatística sobre o limiar.
\end{itemize}

\subsection{Camion em torneios}

Em torneios fortes, o algoritmo de Camion devolve o ciclo em tempo desprezível ($0{,}003$~ms em $n=30$, $29$ nós, um por vértice) enquanto o Held-Karp em $n=14$ já gasta cerca de $0{,}9$~ms e $34\,827$ nós. O mesmo problema, restrito a torneios, é polinomial; isso é o que a seção de aplicações chamou de fronteira da dificuldade.

\section{Conclusão}

O trabalho implementa, com contagem de instruções e medição de tempo, cinco das seis estratégias pedidas (força bruta, backtracking, programação dinâmica, guloso e recozimento como heurística) e justifica por escrito por que divisão e conquista não se sustenta no caso geral. O resultado experimental principal é a separação por ordem de grandeza: força bruta cresce com $(n-1)!$ nós medidos, backtracking diminui a constante e continua fatorial, e Held-Karp cresce com $2^n$, com razão de cerca de $800$ vezes em $n=11$.

\paragraph{Limites.} O teto prático do Held-Karp é a memória, não o tempo: a tabela de predecessores chega a cerca de $385$~MB em $n=22$. O guloso não tem garantia e erra nos contraexemplos da Seção~\ref{sec:guloso}. O recozimento recupera quase todos os ciclos nos testes, mas nunca prova inexistência, e por isso nunca decreta ``não existe'': ele separa ``ciclo válido'' de ``orçamento esgotado''. Os tempos abaixo de $0{,}005$~ms foram descartados por resolução do relógio, e o tipo de limiar teve poucos digrafos com ciclo.

\paragraph{Dificuldades.} (i) A reconstrução do ciclo no Held-Karp exige guardar o predecessor de cada estado, o que multiplica a memória por $n$. (ii) A tentação de apresentar o particionamento em blocos como divisão e conquista; o contraexemplo mostrou o falso negativo, e o trabalho o expõe em vez de o esconder. (iii) O tipo aleatório ingênuo dá digrafos ou triviais (DAG) ou fáceis (arcos extras), e a comparação entre os exatos só ficou informativa com o tipo pior caso. (iv) Distinguir, em todos os métodos, ``não encontrou'' de ``não existe'', para não afirmar o que a heurística não prova.

\begin{thebibliography}{99}
\bibitem{adleman1994} L. M. Adleman. Molecular computation of solutions to combinatorial problems. \emph{Science}, 266(5187):1021--1024, 1994.
\bibitem{bellman1962} R. Bellman. Dynamic programming treatment of the travelling salesman problem. \emph{Journal of the ACM}, 9(1):61--63, 1962.
\bibitem{bjorklund2010} A. Björklund. Determinant sums for undirected Hamiltonicity. In \emph{Proc. 51st IEEE FOCS}, 2010.
\bibitem{camion1959} P. Camion. Chemins et circuits hamiltoniens des graphes complets. \emph{Comptes Rendus de l'Académie des Sciences}, 249:2151--2152, 1959.
\bibitem{compeau2011} P. E. C. Compeau, P. A. Pevzner, G. Tesler. How to apply de Bruijn graphs to genome assembly. \emph{Nature Biotechnology}, 29:987--991, 2011.
\bibitem{cormen2009} T. H. Cormen, C. E. Leiserson, R. L. Rivest, C. Stein. \emph{Introduction to Algorithms}. 3.\ ed., MIT Press/McGraw-Hill, 2009.
\bibitem{frieze1988} A. M. Frieze. An algorithm for finding Hamilton cycles in random directed graphs. \emph{Journal of Algorithms}, 9(2):181--204, 1988.
\bibitem{garey1979} M. R. Garey, D. S. Johnson. \emph{Computers and Intractability: A Guide to the Theory of NP-Completeness}. W. H. Freeman, 1979.
\bibitem{held1962} M. Held, R. M. Karp. A dynamic programming approach to sequencing problems. \emph{Journal of the SIAM}, 10(1):196--210, 1962.
\bibitem{karp1972} R. M. Karp. Reducibility among combinatorial problems. In \emph{Complexity of Computer Computations}, Plenum, 1972, pp.\ 85--103.
\bibitem{kirkpatrick1983} S. Kirkpatrick, C. D. Gelatt, M. P. Vecchi. Optimization by simulated annealing. \emph{Science}, 220(4598):671--680, 1983.
\bibitem{ziviani2006} N. Ziviani. \emph{Projeto de Algoritmos com implementações em Java e C++}. Cengage Learning (Thomson/Pioneira), São Paulo, 1.\ ed., 2006.
\end{thebibliography}

\end{document}
